\section{Basics of Divisibility}
\begin{definition}[Divisibility]
    Let $a,b\in\mathbb{Z}$. We say that $a \neq 0$ divides $b$ if there exists an integer $c$ such that $b=ac$. We write $a\mid b$ if $a$ divides $b$.
\end{definition}
We say $a$ is a divisor of $b$ if $a\mid b$. The set of all divisors of $b$ is denoted by $Div(b)$.
\[
Div(b) := \{a \in \Z: a \mid b\}.
\]
We can also say $a \mid b$ implies there is a solution in $\Z$ to the equation $ax=b$.
\begin{prop}
    \begin{enumerate}
        \item For all $a\in\Z \setminus {0}$, $a\mid 0$.
        \item $a\mid b$ and $b\mid c$ implies $a\mid c$.
        \item $a \mid b$ and $a \mid c$ implies $a \mid (bx+cy)$ for all $x,y\in\Z$.
    \end{enumerate}
\end{prop}
\begin{proof}
    \begin{enumerate}
        \item $a\mid 0$ since $0=a\cdot 0$.
        \item $a\mid b$ implies $b=ax$ and $b\mid c$ implies $c=by=axy=a(xy)$ implies $a \mid c$.
        \item $a\mid b$ implies $b=ax$ and $a\mid c$ implies $c=ay$. So, $bx+cy=a(x^2+y^2)$ implies $a \mid (bx+cy)$.
    \end{enumerate}
\end{proof}
\begin{theorem}[Division Algorithm]
    Let $a,b\in\Z$ with $b>0$. Then there exist unique integers $q,r$ such that $a=qb+r$ and $0\leq r<b$.
\end{theorem}
\begin{proof}
    Consider the set
    \[S = \{a - xb \mid x\in\Z, a-xb \ge 0\}.\]
    The set is nonempty because $a - (-|a|)b = a + |a|b \ge a + |a| \ge 0.$ According to the Well-Ordering Principle, $S$ has a least element, say $r$. So, $r = a - qb$ for some $q\in\Z$. We will show that $0\leq r < b$. $0 \le r$ by the construction of $S$. We argue that $r < b$. If this were not the case, then $r \ge b \implies r - b \ge 0$. But 
    \[
    a - (q+1)b = a - qb - b = r - b \ge 0.
    \]
    Hence, $a - (q+1)b = r - b \in S$ which contradicts the minimality of $r$. Hence, $0\leq r < b$. This completes the proof of existence.

    Now, we will show the uniqueness of $q$ and $r$. Suppose there exist $q_1, q_2, r_1, r_2$ such that $a = q_1b + r_1$, $a = q_2b + r_2$, $0\leq r_1, r_2 < b$. Subtracting the two equations, we get \[(q_1 - q_2)b = r_2 - r_1.\] Since $0\leq r_1, r_2 < b$, we have $-b < r_2 - r_1 < b$, or, \[|r_2 - r_1| < b.\] Therefore, $|q_1-q_2|b < b$, which implies,
    \[0 \le |q_1 - q_2| < 1\]
    So, $r_2 - r_1 = 0 \implies r_1 = r_2$ and $q_1 = q_2$. This completes the proof of uniqueness.
\end{proof}
\section{Primes \& Related Theorems}
\begin{definition}[Prime Number]
    An element $p\in\N$ and $p > 1$ is called a prime number if for all $q \in \N$, $q \mid p \implies q = 1$ or $q = p$.
\end{definition}
Equivalently, $p > 1$ is prime if $Div(p) = \{\pm 1, \pm p\}$.
\begin{eg}
$2, 3, 5, 7, 11, 13, 17, 19, 23, 29, 31, 37, 41, 43, 47, 53, 59$ etc. are primes.
\end{eg}
\begin{eg}[Biggest Prime]
    As of writing this note, the biggest known prime number is \[2^{136279841}-1.\]
\end{eg}
\begin{theorem}[Existence of Prime Factorization]
    Every positive integer greater than $1$ can be written as a product of primes.
\end{theorem}
\begin{proof}
    Assume for the sake of contradiction, that there exits a positive integer greater than $1$ that cannot be written as a product of primes. By well-ordering principle, we will have a least $n$ of this sort. If $n$ is prime, then it is already a product of primes (just itself), which is a contradiction. 

    If $n$ is composite, then it can be written as $n = a \cdot b$, for some $a,b \in \N$ where $1 < a,b < n$. Since $a,b < n$, they can be written as a product of primes. So, $n = a \cdot b$ can also be written as a product of primes, which is again a contradiction.
    
    Hence, every positive integer greater than $1$ can be written as a product of primes.
\end{proof}
We can extend this concept to negative integers as well by including $-1$ as a factor.  
For any integer $n \neq 0$,
\[
n = (-1)^{\epsilon(n)} \prod_{i=1}^{k} p_i^{a(p_i)},
\]
where $\epsilon(n) = 0$ if $n > 0$ and $ 1$ otherwise. And $a(p_i)$ is called \textit{order of $n$ at $p_i$} which is a non-negative integer such that $p^a \mid n$ but $p^{a+1} \nmid n$.

Next lemma requires a theorem which we will just state but not prove.
\begin{theorem}[Bezout's Identity]
    Given two integers $a$ and $b$, not both of which are zero, there exist $x,y\in\Z$ such that \[ax+by=\gcd(a,b).\]
\end{theorem}
\begin{lemma}
    If $p$ is a prime and $p \mid ab$, then $p \mid a$ or $p \mid b$.
\end{lemma}
\begin{proof}
    If $p \mid a$, then we are done. So, assume $p \nmid a$. Let $g = gcd(a, p).$ Since, $p$ is a prime $g = 1$ or $g = p$ (which is not possible). Hence, $g = 1$. By Bezout's Identity, there exist $x,y\in\Z$ such that $px + ay = 1.$ Multiplying both sides by $b$, we get \[bpx + aby = b.\] In the equation, left side is divisible by $p$, so is the right side. Hence, $p \mid b$.
\end{proof}
\begin{corollary}
    If $p$ is a prime and $p \mid a_1a_2\dots a_n$, then $p \mid a_i$ for some $i$.
\end{corollary}
\begin{proof}
    We will prove this by induction on $n$. For $n=1$, the statement is trivial. When $n = 2$, the result is the content of Lemma $2.1$. Assume the result is true for $n = k$. We will show that it is true for $n = k+1$. By Lemma $2.1$, $p \mid a_1a_2\dots a_k$ or $p \mid a_{k+1}$. Let $p \mid a_1a_2\dots a_{k+1}$. By induction hypothesis, $p \mid a_1a_2\dots a_k$ implies $p \mid a_i$ for some $i$ and we are done. Otherwise, $p \mid a_{k+1}$. This completes the proof.
\end{proof}
\begin{theorem}[Fundamental Theorem of Arithmetic]
    Every positive integer greater than $1$ can uniquely be written as a product of primes (upto reordering).
\end{theorem}
\begin{proof}
    Existence is already proved in theorem $2.2$. We will now show uniqueness. Assume for the sake of contradiction that there exists a positive integer greater than $1$ that can be written as a product of primes in two different ways. Let $n$ be the smallest such number. Let $n = p_1p_2\dots p_k = q_1q_2\dots q_l$ where $p_i, q_i$ are primes. Since $p_1 \mid q_1q_2\dots q_l$, by Corollary $2.1$, $p_1 \mid q_i$ for some $i$. WLOG, let $p_1 \mid q_1$. Since $q_1$ is prime, $p_1 = q_1$. Let, $n' = p_2p_3\dots p_k = q_2q_3\dots q_l$. Clearly, $n > n'$ which contradicts the minimality of $n$. Hence, every positive integer greater than $1$ can uniquely be written as a product of primes.
\end{proof}
\begin{theorem}
    There are infinitely many primes.
\end{theorem}
\begin{proof}
    Assume for the sake of contradiction that there are finitely many primes, say $p_1, p_2, \dots, p_n$. Consider the number $N = p_1p_2\dots p_n + 1$. Since $N > 1$, it can be written as a product of primes. But none of $p_1, p_2, \dots, p_n$ divides $N$ as $N$ leaves a remainder of $1$ when divided by any of $p_i$. Hence, $N$ is a prime which is not in the list $p_1, p_2, \dots, p_n$. This contradicts the assumption that there are finitely many primes.
\end{proof}
\begin{definition}[Riemann Zeta Function]
    The Riemann Zeta function is defined as
    \[\zeta(s) = \sum_{n=1}^{\infty} \frac{1}{n^s},\]
    where the function is defined for $s \in \R$ (in this context) and $s > 1$.
\end{definition}
\begin{theorem}[Euler's Product Formula]
    For $s > 1$, we have
    \[\zeta(s) = \prod_{p \text{ prime}} \frac{1}{1-p^{-s}}.\]
\end{theorem}
\begin{proof}
    The geometric series formula gives us,
    \[\frac{1}{1-p^{-s}} = \sum_{m=0}^{\infty} \frac{1}{p^{ms}}.\]
    So,
    \[\prod_{p \text{ prime}} \frac{1}{1-p^{-s}} = \prod_{p \text{ prime}} \sum_{m=0}^{\infty} \frac{1}{p^{ms}} = \sum_{n=1}^{\infty} a_n \frac{1}{n^s}.\]
    By Fundamental Theorem of Arithmetic, $a_n = 1$ for all $n$. Hence, the product is equal to $\zeta(s)$.
\end{proof}
Now, set $s=1$ in Euler's Product Formula to get 
\[\zeta(1) = \sum_{n=1}^{\infty} \frac{1}{n} = \prod_{p \text{ prime}} \frac{1}{1-p^{-1}}.\]
Left side is the harmonic series which diverges to $\infty$. So, the product on the right side must also diverge to $\infty$ which requires the product on the right side to be over infinite index. This implies that we must have infinitely many primes.
\section{Distribution of Primes}
Gauss conjectured that the density of primes is approximately given by: 
\[
\frac{1}{\log x}
\]
Let \( \pi(x) \) denote the number of primes less than or equal to \( x \). Then, an asymptotic estimate for \( \pi(x) \) is: 
\[
\pi(x) \sim \frac{x}{\log x}
\]
In asymptotic notation, if \( f(x) \sim g(x) \), it means: 
\[
\lim_{x \to \infty} \frac{f(x)}{g(x)} = 1.
\]
A more refined approximation is given by the logarithmic integral: 
\[
\text{Li}(x) = \int_2^x \frac{dt}{\log t}.
\]
The \textbf{Prime Number Theorem (PNT)} states that: 
\[
\pi(x) \sim \text{Li}(x),
\]
which was proven in 1896.

\textbf{Von Mangoldt's Explicit Formula:} Von Mangoldt’s explicit formula connects the sum of the von Mangoldt function $\Lambda(n)$ to the nontrivial zeros of the Riemann zeta function. It is a key result in analytic number theory, particularly in understanding the distribution of prime numbers.

Define the von Mangoldt function \( \Lambda(n) \) as:
\[
\Lambda(n) =
\begin{cases}
\log p, & \text{if } n = p^m \text{ for some prime } p \text{ and integer } m \geq 1, \\
0, & \text{otherwise}.
\end{cases}
\]
The Chebyshev function \( \psi(x) \) is given by:
\[
\psi(x) = \sum_{n \leq x} \Lambda(n).
\]
Von Mangoldt's explicit formula states:
\[
\psi(x) = x - \sum_{\rho} \frac{x^{\rho}}{\rho} - \log 2\pi - \frac{1}{2} \log(1 - x^{-2}),
\]
where the sum runs over all the nontrivial zeros \( \rho \) of the Riemann zeta function \( \zeta(s) \).